\section{System Architecture}

Figure~\ref{fig:overview} shows the end-to-end architecture: Eventsim generates events into a Kafka cluster, Spark Structured Streaming consumes and lands them as Hudi bronze tables, dbt transforms the bronze layer into a star schema and a wide table on the same Spark cluster, and Superset visualizes the result through the Spark Thrift Server, alongside a Prometheus/Grafana monitoring stack.

\begin{figure}[H]
    \centering
    \includegraphics[width=0.95\textwidth]{overview.png}
    \caption{Streamify architecture overview: Eventsim, Kafka, Spark Structured Streaming, Hudi, dbt, Superset, and monitoring.}
    \label{fig:overview}
\end{figure}

\begin{table}[H]
    \centering
    \caption{Technology stack per pipeline layer.}
    \label{tab:stack}
    \begin{tabularx}{\textwidth}{|l|l|X|}
        \hline
        \textbf{Layer} & \textbf{Tool} & \textbf{Role} \\
        \hline
        Data simulation & Eventsim & Generates synthetic music-streaming events based on the Million Song Dataset \\
        \hline
        Message queue & Apache Kafka 4.3 (KRaft) & 2-node cluster, replication factor 2, 3 partitions per topic, no ZooKeeper \\
        \hline
        Stream processing & Spark Structured Streaming 3.5.3 & Reads Kafka, decodes/enriches events, writes bronze tables to Hudi \\
        \hline
        Table format / storage & Apache Hudi & MERGE\_ON\_READ bronze streaming tables, COPY\_ON\_WRITE gold/seed tables, on MinIO/S3 via the Hadoop S3A connector \\
        \hline
        Compute cluster & Spark Standalone (1 master, 2 workers) & Shared by the streaming job and dbt (via Thrift), 4 cores total \\
        \hline
        Transformation & dbt (\texttt{dbt-spark}, \texttt{method: thrift}) & SQL-based star schema and wide-table modelling, submitted via the Spark Thrift Server \\
        \hline
        Visualization & Apache Superset & Dashboards connected to the Spark Thrift Server \\
        \hline
        Monitoring & Prometheus, Grafana, cAdvisor & Kafka consumer-group lag, Spark JVM/streaming metrics, per-container CPU/memory \\
        \hline
        Containerization & Docker Compose & Single-command reproducible deployment \\
        \hline
    \end{tabularx}
\end{table}

\section{Implementation Workflow}

The pipeline runs as the following seven stages:

\begin{enumerate}
    \item \textbf{Data simulation (Eventsim).} Eventsim simulates a real-world user base, generating page-view and song-play events.
    \item \textbf{Ingestion and messaging (Kafka).} Apache Kafka 4.3 runs in KRaft mode (no ZooKeeper) as a 2-node cluster. Each topic has 3 partitions with replication factor 2 for fault tolerance.
    \item \textbf{Real-time processing (Spark Structured Streaming).} A long-running Spark job consumes all three topics, fixes latin1-mojibake on text fields, derives time columns, and appends each micro-batch to bronze Hudi tables (\texttt{staging.<topic>}), using the MERGE\_ON\_READ table type so streaming appends land as delta log files that are periodically compacted into columnar base files.
    \item \textbf{Lakehouse storage (Hudi).} Tables are Hudi-managed, backed by MinIO/S3 locally or real S3 in \texttt{aws} mode via the Hadoop S3A connector, with no rebuild needed to switch. There is no separate catalog service: each table is self-describing at its storage path (a \texttt{.hoodie/} metadata directory) and is (re-)registered into Spark's catalog idempotently at job startup.
    \item \textbf{Data transformation (dbt).} dbt~\cite{dbt} connects to the Spark cluster via the Spark Thrift Server (\texttt{method: thrift}) and builds a star schema plus a denormalized wide table directly on top of the bronze Hudi tables, materialized as COPY\_ON\_WRITE Hudi tables rebuilt via full overwrite each run.
    \item \textbf{Business intelligence (Superset).} Apache Superset~\cite{superset} connects to the same Spark Thrift Server, querying the \texttt{wide\_streams} model to visualize insights through interactive dashboards.
    \item \textbf{Monitoring (Prometheus and Grafana).} Prometheus scrapes Kafka consumer-group lag (via kafka-exporter), Spark's built-in Prometheus metrics, and per-container stats (via cAdvisor); Grafana visualizes all of it across dashboards.
\end{enumerate}

\subsection{Star Schema and Wide Table}

dbt builds two model layers on top of the bronze Hudi tables:

\begin{itemize}
    \item \textbf{Star schema.} One fact table, \texttt{fact\_streams}, joined to five dimension tables: \texttt{dim\_users}, \texttt{dim\_songs}, \texttt{dim\_artists}, \texttt{dim\_locations}, and \texttt{dim\_datetime}. This is the normalized, storage-efficient model (Section~\ref{sec:core-theory}).
    \item \textbf{Wide table (\texttt{wide\_streams}).} A denormalized table built from the star schema by pre-joining \texttt{fact\_streams} with all five dimension tables, so that Superset can query a single flat table instead of performing joins at dashboard-query time.
\end{itemize}

\begin{figure}[H]
    \centering
    \includegraphics[width=0.7\textwidth]{star_schema.png}
    \caption{Star schema: \texttt{fact\_streams} and its five dimension tables.}
    \label{fig:star-schema}
\end{figure}

\begin{figure}[H]
    \centering
    \includegraphics[width=0.7\textwidth]{dbt_graph.png}
    \caption{dbt DAG showing the dependency between bronze sources, the star schema, and \texttt{wide\_streams}.}
    \label{fig:dbt-graph}
\end{figure}

\subsection{Streaming Job: Kafka to Hudi Bronze}

Listing~\ref{lst:build-stream} shows the core of the Spark Structured Streaming job that reads a Kafka topic, decodes it, and derives the time columns used to partition the bronze Hudi tables.

\begin{lstlisting}[language=Python, caption={Building a Hudi bronze stream from a Kafka topic (excerpt from \texttt{spark\_streaming/jobs/stream\_events.py}).}, label={lst:build-stream}]
def build_stream(spark, topic, decode_strings):
    raw = (
        spark.readStream.format("kafka")
        .option("kafka.bootstrap.servers", KAFKA_BROKERS)
        .option("subscribe", topic)
        .option("startingOffsets", "earliest")
        .load()
    )

    parsed = raw.select(
        F.from_json(F.col("value").cast("string"), SCHEMAS[topic]).alias("data"),
        "partition",
        "offset",
    ).select("data.*", "partition", "offset")

    if decode_strings:
        parsed = (
            parsed.withColumn("artist", decode_udf(F.col("artist")))
                  .withColumn("song", decode_udf(F.col("song")))
        )

    ts_col = F.timestamp_millis(F.col("ts"))
    return (
        parsed
        .withColumn("ts_timestamp", ts_col)
        .withColumn("year_", F.year(ts_col))
        .withColumn("month_", F.month(ts_col))
        .withColumn("day_", F.dayofmonth(ts_col))
        .withColumn("hour_", F.hour(ts_col))
        .withColumn("processed_time", F.current_timestamp())
        .withColumn("_row_key", F.expr("uuid()"))
    )
\end{lstlisting}

\section{Cluster and Environment Setup}

The whole stack is deployed with Docker Compose (\texttt{docker-compose.yml}), with a \texttt{docker-compose.dev.yml} overlay that adds MinIO for local, S3-free development. Spark runs as a Standalone cluster of 1 master and 2 workers, 2 cores / 2 GB each (4 cores total), shared between the streaming ingestion job and the Spark Thrift Server that dbt submits SQL through, each capped at 2 cores.
